\chapter*{概要}
\addcontentsline{toc}{chapter}{概要}

本研究は，個人の性格特性と日々の推し活の状況から，その人に適した推し活方法を提案するシステム「OshiFit（オシフィット）」を構築し，その有用性を検証することを目的とする．

\paragraph{研究背景}
近年，精神的ストレスによるメンタル不調を抱える人は増加しており，精神疾患を有する総患者数は平成29年時点で約419.3万人に達している．ストレスを和らげる手段として，睡眠や運動に加えて趣味活動が注目されており，その中でも特定のアイドル，キャラクター，配信者などを継続的に応援する「推し活」は，Z世代を中心に広く行われている．推し活は「生きがい」や自己表現，人とのつながりをもたらす一方で，睡眠不足，過度な出費，「追わなければならない」という義務感，SNS上での他のファンとの比較による落ち込みといった負担（本研究では「推し疲れ」と呼ぶ）も生む．スクリーンタイムの制限や家計簿アプリは「減らすほど良い」ことを前提とするが，推し活では負担の原因となる活動が同時に生活の張りの源でもあるため，一律に減らす助言は楽しさまで奪う可能性がある．さらに，本人は自分のどの習慣が推し疲れにつながっているのかに気づきにくい．先行研究の多くは推し活の動機や心理を記述するにとどまり，便益と負担を同時に扱った定量的な検討や，個人の状態に合わせて推し活の仕方を支援する仕組みは十分に検討されていない．筆者の前研究（卒業論文）では，推し活時の心拍とストレスの関連を体験者3名・2週間で調べたが，体験者が少なく，受動的な推し活に限られ，期間も短いという課題が残った．

\paragraph{研究目的とリサーチクエスチョン}
本研究では，多くのファンの性格特性・推し活行動・幸福度・推し疲れの関係を分析し，その知見に基づいて，本人が自ら選べる提案と「今のまま続けた場合（プランA：現在のままの推し活）」と「提案を取り入れた場合（プランB：最適化された推し活）」の28日後の見込みを，自分に似せたアバターと数値で示すシステムを構築する．本人が気づきにくい負担の原因に気づき，自ら選んだ小さな変更によって幸福度の向上とストレスの軽減につなげることをねらう．次の3つのリサーチクエスチョン（RQ）を設定した．RQ1：性格特性（TIPI-J）と推し活行動は，幸福度・推し疲れとどのように関連し，個人の見込みをどの程度の精度で予測できるか．RQ2：約28日間の利用で，すでに良好な人（群A），プランBを選んだ人（群B），プランAを選んだ人（群C）の幸福度と推し疲れはそれぞれどう変化し，群の間でどう異なるか．RQ3：利用者は自分の推し疲れの原因に気づき，特に行動を変えた人はシステムを有用と評価するか．

\paragraph{調査（RQ1）}
推し活を行う日本のファンを対象にオンライン質問紙調査を行い，回答者234名から，推し活をしていない17名と注意確認項目に誤答した2名を除く215名（男性113名，女性94名，その他・無回答8名，年齢の中央値19歳）を分析した．性格特性はTIPI-Jの10項目（5件法），推し活行動は視聴・現地イベント・グッズ購入・SNSなどの頻度，推し活への没入，推し活をしたくなるきっかけ，推し活による幸福度と推し疲れを測定した．k-meansによる類型化，観測変数によるパスモデル，交互作用の検定，10分割交差検証による予測精度の推定を行った．その結果，79\%が推しを「生きがい」と答える一方で，40\%がお金の不安，33\%が睡眠への支障を報告し，幸福度と推し疲れは正の相関（$r=.24$）を示した．推し活の頻度は幸福度（$\beta=.43$）と推し疲れ（$\beta=.29$）の両方を高め，没入も両方に関連した（$\beta=.20$，$.23$）．すなわち，便益と負担は同じ行動から生じていた．性格特性は推し活の関わり方と関連した（外向性→頻度 $\beta=.19$，誠実性→没入 $\beta=-.15$，神経症傾向→気分の回復のための推し活 $\beta=.21$，開放性→同 $\beta=-.16$）が，頻度の効果を調整する効果は認められなかった（交互作用を加えても決定係数はほぼ変わらず，$p=.74$，$.79$）．性格のみによる類型化は分離が弱く（シルエット係数0.15〜0.18），性格は単独の「タイプ」としてではなく，行動と組み合わせて用いるのが適切であることが示された．性格・行動・きっかけを組み合わせた予測の交差検証決定係数は幸福度で.27，推し疲れで.11にとどまり，個人の予測には大きな不確かさが伴うことが分かった．

\paragraph{システム（OshiFit）}
これらの結果から，(D1)一律に減らす助言をせず「今のまま」も常に有効な選択肢とする，(D2)固定的な性格タイプではなく性格と本人の行動・負担を組み合わせて個人化する，(D3)便益と負担を同時に示す，(D4)予測の不確かさを正直に示す，という4つの設計指針を導いた．OshiFitは，質問紙（TIPI-J，推し活の内容の割合，気持ちと生活の変化，推し活をしたくなるとき，WHO-5，SHS，PSS-J．検証実験では短縮版の68問）から，幸福度・推し活の充実・睡眠と生活・お金・義務感と人間関係・SNSと比較の6軸を0〜100点で算出し，算出式と使用項目を画面に表示する．6軸のうち5軸以上が60点以上かつWHO-5の素点が13点以上の人には「今のままで順調」と伝え（群A），それ以外の人には，点数の低い軸に対応する15種類の提案から，性格の傾向と活動の比重も考慮して最大3案を示す．さらに，プランAとプランBの28日後の見込みを80\%の幅とともに示し，利用者が部品を選んで作った本人に似せたアバター（写真は使わない）の表情で，見込みの差を表す．利用者はプランBを選ぶ（群B）か，プランAを選ぶ（群C）かを自分で決め，約28日間，毎日のアンケート（約1分）と4日ごとの振り返り（全6回）で提案の実行状況と気分・推し疲れを記録する．見込みは本人の記録で少しずつ補正する．

\paragraph{検証実験（RQ2，RQ3）}
大阪工業大学の研究倫理審査の承認を得て，約28日間の検証実験を行う．群は無作為ではなく，開始時の状態（群A）と本人の選択（群B・C）で決まる．これは，趣味である推し活に変更を強制することが設計指針D1に反するためである．主要な結果指標はWHO-5と推し疲れ（総合）であり，各群の開始時と最後の調査の比較，開始時の値を共変量とした群Bと群Cの共分散分析，傾向スコアによる重み付けでの感度分析，毎日の記録の線形混合モデル，提案の実行の程度と変化量の関係（用量反応），見込みの当たり方を分析する．\todo{実験後に記入：参加者数，各群の人数，WHO-5と推し疲れの変化量と検定結果，気づき・有用性の評価結果}

\paragraph{本研究の意義}
本研究の意義は，第一に，推し活の便益と負担が同じ行動から生じることを215名の調査で定量的に示し，「減らせば良い」という支援が推し活には適さないことを明らかにした点にある．第二に，性格特性を固定的なタイプとしてではなく，行動・負担と組み合わせて個人化に用いるという設計の根拠を示した点にある．第三に，予測の不確かさを隠さずに示しながら，本人に似たアバターで2つの未来を比較させ，本人が自ら選ぶ行動変容を支援するシステムを構築した点にある．\todo{実験後に追記：検証実験から得られた知見}本研究の知見は，推し活に限らず，楽しさと負担が表裏一体となった趣味や娯楽を，やめさせるのではなく持続可能な形に整えるための支援の設計に応用できると考えられる．

\vspace{1em}
\noindent\textbf{キーワード：}推し活，ウェルビーイング，推し疲れ，性格特性（TIPI-J），パーソナルインフォマティクス，アバター，予測の不確かさ
